
% ФПВ берется из ЛВ

\color{unimain}{\subsection{Фиксация положения выключателей (ФПВ)}}
\color{black}


\begin{enumerate}

\item
Схемы логики ФПВ В2(ОВ)ВН, ФПВ В3, ФПВ В1(ЛВ)СН, ФПВ В2(ОВ)СН аналогичны логике ФПВ В1(ЛВ)ВН.

\item
Схемы логики Р2, Р3 аналогичны логике Р1.

\item
Ввод в работу ФПВ В1(ЛВ)ВН осуществляется с помощью уставки \textbf{«ФПВ В1(ЛВ)ВН»}. Оперативный ввод/вывод ФПВ производится с помощью БУ (блоков управления функциями) «Вывод ФПВ В1(ЛВ)ВН». Предусматривается возможность формирования положения выключателей с контролем разъединителей, которая задается уставкой \textbf{«Контроль Р В1(ЛВ)ВН»}. С помощью уставки \textbf{«Схема Р В1(ЛВ) ВН»} имеется возможность выбрать схему выключателя, два или три разъединителя. С помощью уставки \textbf{«Тип В1(ЛВ) ВН»} имеется возможность выбрать тип выключателя, линейный или обходной. Так же от блок-контактов В1(ЛВ) ВН формируется внутренний пуск ЗНФ, который пофазно контролирует расхождение контактов полюсов выключателей, возникающее при трехфазном включении или отключении выключателя.

\begin{figure}[H]
\centering
\includegraphics[width=0.6\textwidth,height=0.6\textheight,keepaspectratio]{\fbpath/06.02 ДЗАТ/761.3471 ФПВ/img/FPV_1.pdf}
%\captionof{figure}{Функциональная схема ФПВ В1}\label{pic:FPV_1}
\end{figure}

\begin{figure}[H]
\centering
\includegraphics[width=0.6\textwidth,height=0.6\textheight,keepaspectratio]{\fbpath/06.02 ДЗАТ/761.3471 ФПВ/img/FPV_2.pdf}
\captionof{figure}{Функциональная схема ФПВ В1}\label{pic:FPV_2}
\end{figure}


\end{enumerate}